\documentclass[12pt]{article}

\usepackage{array}
\usepackage{float}
\usepackage{xcolor}
\usepackage{authblk}
\usepackage{comment}
\usepackage{parskip}
\usepackage{tabularx}
\usepackage{geometry}
\usepackage{setspace}
\usepackage{graphicx}
\usepackage{hyperref}
\usepackage{microtype}
\usepackage[T1]{fontenc}
\usepackage[utf8]{inputenc}
\usepackage{amsmath, amssymb, amsfonts}

\setstretch{1.5}
\geometry{paper=letterpaper, layout=letterpaper, margin=1in}

\begin{comment}

    Package Description

    float       ->      Improves control over the placement of figures and tables.
    xcolor      ->      Adds support for colored text, backgrounds, and page elements.
    authblk     ->      Provides support for author blocks and affiliations.
    fontenc     ->      Selects font encodings for better character output and hyphenation.
    amsmath     ->      Provides enhanced environments and commands for mathematical typesetting.
    comment     ->      Allows block of text to be excluded from the document.
    parskip     ->      Adjusts paragraph spacing and indentation.
    amssymb     ->      Provides additional mathematical symbols.
    geometry    ->      Sets the paper size and layout, including margins.
    inputenc    ->      Specifies the input text encoding, such as UTF-8.
    graphicx    ->      Enables including and scaling images.
    setspace    ->      Controls line spacing, such as single or double spacing.
    amsfonts    ->      Provides additional mathematical fonts, including \mathbb.
    microtype   ->      Improves typography through subtle spacing and character adjustments.

\end{comment}

\begin{document}

\title{Relationship between the 3D Printing Parameters and the Nonlinear Properties of Auxetic Metamaterials, and Development of A Mathematical Model Using the Neural Network to Predict This Relationship}

% Please confirm your full name here. Thanks a lot!

% Keep the 8th position empty, here, the advisor or
% supervisor name will be confirmed.

\author[1]{Abdullah al Azmi}
\author[2]{Habibunnabi Nahid}

\affil[1-2]{
    Department of Mechanical and Production Engineering
    \protect \\
    Islamic University of Technology (IUT), Gazipur, Bangladesh
}

\date{\today}

\maketitle

\begin{abstract}

    3D-printed auxetic metamaterials have gained significant interest in recent years due to the ease of fast prototyping. While previous studies examine the nonlinear properties of the auxetic metamaterials from the geometrical perspective, including unit cell shape, direction, type, and magnitude, none of the studies closely examine how the nonlinear properties change with 3D printing parameters, including infill patterns, layer height, layer thickness, infill density, and so on. Hence, this study aims to closely examine how the nonlinear properties of the auxetic metamaterials vary with 3D printing parameters and to develop a mathematical model to predict these relationships using artificial neural networks and to inversely design the auxetic metamaterials.

\end{abstract}
\section{Introduction}
\input{Sections/related_work.tex}
\section{Methodology}
\section{Conclusion}
\section{References}

\end{document}
